\section{Introduction}
\label{sec:introduction}

Conformer generators provide candidate three-dimensional geometries for
molecular modeling when the required structure is not known in advance.
Molecular shape comparison and pharmacophore screening use these geometries
to compare spatial arrangements of atoms or interaction features, while
docking searches for ligand geometries compatible with a binding site.
Even docking methods that allow ligand flexibility may sample only some
internal degrees of freedom, making their results dependent on the supplied
conformations.\cite{mcnutt2023conformers} Experimentally determined
protein--ligand structures provide reference geometries for assessing whether
a generated ensemble contains a bioactive conformation.

Previous work has established that ensemble construction affects the usefulness
of generated conformations. McNutt et al.\ examined RDKit ETKDG and the deep
generative model DMCG, varying ensemble size, conformer selection, and energy
minimization, and evaluated bioactive conformation recovery, pharmacophore
screening, and docking.\cite{mcnutt2023conformers} That study connected
generation choices to application performance using a focused comparison of
two generators. Selecting among the wider range of available generators
requires a complementary comparison of their ability to recover observed
bound geometries while preserving molecular identity and local geometric
plausibility under specified sampling allowances.

The available methods span substantially different sampling procedures.
Distance geometry combines molecular connectivity with geometric constraints
and experimental torsional preferences.\cite{wang2020etkdg} Learned
approaches include diffusion over torsions, continuous conformer fields,
and molecular language representations.\cite{jing2022torsional,wang2024mcf,liu2025nextmol}
LoQI and its associated ChEMBL3D ensembles emphasize low-energy
structures,\cite{nikitin2025loqi} while FlowR provides another learned
approach to coordinate generation.\cite{cremer2025flowr} Their different
representations, refinement procedures, and filtering policies make it
necessary to evaluate the resulting ensembles against common structural
criteria. Performance of one learned generator cannot establish the
suitability of these other methods for recovering bioactive conformations.

Here, we define an evaluation protocol for selecting generators of ensembles
intended to recover bioactive conformations. A successful ensemble must
contain at least one conformer that preserves the requested molecular
identity, passes intramolecular plausibility checks, and lies within
0.75~\AA{} heavy-atom root mean square deviation (RMSD) of the experimental
geometry after alignment. PoseBusters supplies the identity and plausibility
checks.\cite{buttenschoen2024posebusters} We compare two candidate allowances:
the stored ChEMBL3D count for each ligand and 1,000 conformers per ligand.
These allowances apply before filtering; retained counts are reported
alongside recovery. Geometric diversity is evaluated as a possible indicator
of recovery, rather than imposed as a minimum requirement: a diverse pool
need not contain a close match to the bound geometry.

Our main comparison includes four classical generation pipelines and six
learned generators: LoQI, Torsional Diffusion, Molecular Conformer Fields,
NExT-Mol, FlowR, and an autoregressive Qwen model. This broader method panel
extends the focused RDKit--DMCG comparison to assess which evaluated
generators satisfy the same recovery criterion under the specified validity
and sampling conditions. The primary cohort comprises 94 CASF-2016 ligand
entries matched to ChEMBL3D; its stored ensembles provide an additional
comparison.\cite{su2019casf,nikitin2025loqi} Analyses by molecular size and
flexibility identify where each generator fails to recover the bound
geometry. Results for a larger reference cohort are provided in the
Supporting Information.

Recovery of one bound structure does not describe how often a generator
samples near observed conformations or whether it covers other observed
geometries of the same molecule. A separate exploratory panel of 23 druglike
molecules with multiple experimental references distinguishes reference
coverage from the fraction of generated samples close to those references.
Together, these evaluations assess recovery and sample concentration for
generator selection. The multi-reference comparison remains subject to
differences in filtering and failure handling described in Methods.
The study evaluates structural recovery rather than downstream screening or
docking performance, and the experimental references are not assumed to
exhaust the accessible conformations of a molecule.
% AUTHOR: resolve dataset provenance, checkpoint selection and training overlap
% audits before claiming performance on unseen molecules or broad generalization.
